% ============================================================
% Mise en page des chapitres de cours compilés (content/cours/<id>/cours.tex).
% À charger APRÈS preamble.tex : reprend le format du recueil COMIMa
% « Techniques de résolution de problèmes » (A5, Latin Modern, en-tête
% « Chapitre N. Titre » / numéro de page).
% Paquets hors preamble.tex : mdframed, needspace, etoolbox (texlive-latex-extra).
% ============================================================
\geometry{a5paper,top=1.8cm,bottom=1.6cm,left=1.5cm,right=1.5cm,headsep=0.5cm}
\usepackage[framemethod=default]{mdframed}
\usepackage{needspace}
\usepackage{etoolbox}

% Pas de titre de section isolé en bas de page.
\pretocmd{\section}{\needspace{8\baselineskip}}{}{}
\pretocmd{\subsection}{\needspace{6\baselineskip}}{}{}

% ---------- Chapitre ----------
% \chapitre{numéro}{titre} : grand titre de début de chapitre, fixe aussi le
% numéro utilisé par les sections (8.1, 8.2…) et l'en-tête courant.
\newcounter{chapitre}
\newcommand{\titrechapitre}{}
\newcommand{\chapitre}[2]{%
  \setcounter{chapitre}{#1}%
  \renewcommand{\titrechapitre}{#2}%
  \thispagestyle{premierepage}%
  \vspace*{1.2cm}%
  {\raggedright\large\bfseries\MakeUppercase{Chapitre}~#1\par}%
  \vspace{0.5cm}%
  {\raggedright\Huge\bfseries #2\par}%
  \vspace{0.8cm}%
}
\renewcommand{\thesection}{\thechapitre.\arabic{section}}

% ---------- En-têtes ----------
\fancyhf{}
\renewcommand{\headrulewidth}{0pt}
\fancyhead[L]{\small Chapitre~\thechapitre. \titrechapitre}
\fancyhead[R]{\small\thepage}
\fancypagestyle{premierepage}{%
  \fancyhf{}%
  \fancyhead[R]{\small\thepage}%
}

% ---------- Énoncés mathématiques ----------
% Numérotation par chapitre : Théorème 8.1, Exemple 8.2…
\theoremstyle{plain}
\newtheorem{theoreme}{Théorème}[chapitre]
\newtheorem{proposition}[theoreme]{Proposition}
\theoremstyle{definition}
\newtheorem{definition}[theoreme]{Définition}
\newtheorem{exemple}[theoreme]{Exemple}
\theoremstyle{remark}
\newtheorem*{remarque}{Remarque}

% Encadré à filet gauche (sécable) pour théorèmes et définitions.
\mdfdefinestyle{filet}{%
  topline=false,bottomline=false,rightline=false,
  linewidth=1.6pt,linecolor=comimagreen,
  innerleftmargin=8pt,innerrightmargin=0pt,
  innertopmargin=2pt,innerbottommargin=2pt,
  leftmargin=0pt,rightmargin=0pt,
  skipabove=8pt,skipbelow=6pt,
}
\surroundwithmdframed[style=filet]{theoreme}
\surroundwithmdframed[style=filet]{proposition}
\surroundwithmdframed[style=filet]{definition}

% Encadré « À retenir » : filet rouge.
\newenvironment{retenir}
  {\begin{mdframed}[style=filet,linecolor=comimared]}
  {\end{mdframed}}

% ---------- Exercices guidés et problèmes ----------
% Les énoncés et solutions sont ceux de content/exercices/<id>/ (inclus,
% jamais recopiés). \exercice et \probleme prennent l'id puis le titre ;
% \solutionde{id}{titre} affiche la solution correspondante en fin de
% chapitre (numéro retrouvé par \ref).
\newcounter{exercice}
\newcounter{probleme}
\newif\ifsolutions

\newcommand{\exercice}[2]{%
  \par\needspace{6\baselineskip}\medskip
  \refstepcounter{exercice}\label{exo:#1}%
  \noindent{\bfseries Exercice guidé~\theexercice\ --- #2.}\par\nopagebreak
  \input{#1/enonce.tex}\par\medskip
}
\newcommand{\probleme}[2]{%
  \par\needspace{6\baselineskip}\medskip
  \refstepcounter{probleme}\label{pb:#1}%
  \noindent{\bfseries Problème~\theprobleme\ --- #2.}\par\nopagebreak
  \input{#1/enonce.tex}\par\medskip
}
% #3 : « exo » ou « pb » selon la liste d'origine.
\newcommand{\solutionde}[3]{%
  \par\needspace{6\baselineskip}\medskip
  \noindent{\bfseries\color{comimagreen}%
    \csname solprefix@#3\endcsname~\ref{#3:#1} --- #2.}\par\nopagebreak
  \input{#1/solution.tex}\par\medskip
}
\expandafter\def\csname solprefix@exo\endcsname{Exercice guidé}
\expandafter\def\csname solprefix@pb\endcsname{Problème}

\setlength{\parindent}{0pt}
\setlength{\parskip}{3pt plus 1pt}
